\chapter{Übersetzungen}
\label{ch:quality/translations}

Wenn deine Mods und Spiele übersetzt werden können, haben mehr Menschen
Freude daran. Laut Google Play haben 64 \% der Luanti-Nutzer unter
Android nicht Englisch als Hauptsprache. Für alle Plattformen zusammen
erhebt Luanti keine Statistik zu den Sprachen der Nutzer, aber der
Anteil nicht englischsprachiger Spieler dürfte hoch sein.

In Luanti schreibst du deine Texte auf Englisch und bildest sie mit
Übersetzungsdateien auf andere Sprachen ab. Übersetzt wird auf dem
Client jedes Spielers, sodass jeder Spieler eine andere Sprache sehen
kann.

\section{Wie funktioniert die Übersetzung auf dem
Client?}\label{wie-funktioniert-die-uxfcbersetzung-auf-dem-client}

\subsection{Markierter Text}\label{markierter-text}

Der Server muss den Clients mitteilen, wie Text übersetzt werden soll.
Dazu werden Steuerzeichen in den Text eingefügt, die Luanti sagen, wo
und wie übersetzt werden soll. Das nennt man \emph{markierten Text}
(\emph{marked up text}); dazu später mehr.

Um Text als übersetzbar zu markieren, verwendest du eine
Übersetzerfunktion (\texttt{S()}), die du mit
\texttt{core.\allowbreak{}get\_\allowbreak{}translator(textdomain)}\index{core@\texttt{core}!get_translator@\texttt{core.\allowbreak{}get\_\allowbreak{}translator}}
erhältst:

\begin{Shaded}
\begin{Highlighting}[]
\KeywordTok{local} \ConstantTok{S} \OperatorTok{=} \VariableTok{core}\OperatorTok{.}\NormalTok{get\_translator}\OperatorTok{(}\StringTok{"mymod"}\OperatorTok{)}

\VariableTok{core}\OperatorTok{.}\NormalTok{register\_craftitem}\OperatorTok{(}\StringTok{"mymod:item"}\OperatorTok{,} \OperatorTok{\{}
    \VariableTok{description} \OperatorTok{=}\NormalTok{ S}\OperatorTok{(}\StringTok{"My Item"}\OperatorTok{),}
\OperatorTok{\})}
\end{Highlighting}
\end{Shaded}

Das erste Argument von \texttt{get\_\allowbreak{}translator} ist die
\emph{Textdomain\index{textdomain@Textdomain}}, eine Art
Namensraum.\index{namensraum@Namensraum} Statt alle Übersetzungen einer
Sprache in einer Datei zu sammeln, werden sie nach Textdomains getrennt,
mit einer Datei je Textdomain und Sprache. Die Textdomain sollte dem
Mod-Namen entsprechen, um Konflikte zwischen Mods zu vermeiden.

Markierten Text kannst du fast überall verwenden, wo lesbarer Text
erwartet wird, etwa in Formspecs,\index{formspec@Formspec} Feldern von
Item-Definitionen, \texttt{infotext} und vielem mehr. In Formspecs musst
du markierten Text mit
\texttt{core.\allowbreak{}formspec\_\allowbreak{}escape}\index{core@\texttt{core}!formspec_escape@\texttt{core.\allowbreak{}formspec\_\allowbreak{}escape}}
maskieren.

Trifft der Client auf übersetzbaren Text, etwa in \texttt{description},
sucht er ihn in der Übersetzungsdatei für die Sprache des Spielers.
Findet er keine Übersetzung, verwendet er den englischen Text.

Übersetzbarer markierter Text enthält den englischen Ausgangstext, die
Textdomain und alle weiteren Argumente von \texttt{S()}. Er ist im
Grunde eine Textkodierung des \texttt{S}-Aufrufs mit allen nötigen
Informationen.

Eine weitere Art markierten Texts liefert
\texttt{core.\allowbreak{}colorize}\index{core@\texttt{core}!colorize@\texttt{core.\allowbreak{}colorize}}.

\begin{hinweis}{Kodierung von markiertem Text}

Das genaue Format von markiertem Text musst du nicht kennen, es hilft
aber beim Verständnis.

\begin{Shaded}
\begin{Highlighting}[]
\NormalTok{"\textbackslash{}27(T@mymod)Hello everyone!\textbackslash{}27E"}
\end{Highlighting}
\end{Shaded}

\begin{itemize}
\tightlist
\item
  \texttt{\textbackslash{}27} ist das Escape-Zeichen. Es sagt Luanti,
  dass etwas Besonderes folgt. Es wird sowohl für Übersetzungen als auch
  für das Einfärben von Text verwendet.
\item
  \texttt{(T@mymod)} besagt, dass der folgende Text mit der Textdomain
  \texttt{mymod} übersetzbar ist.
\item
  \texttt{Hello everyone!} ist der übersetzbare englische Text, wie er
  der Übersetzerfunktion übergeben wurde.
\item
  \texttt{\textbackslash{}27E} ist wieder das Escape-Zeichen, gefolgt
  von \texttt{E}; es markiert das Ende.
\end{itemize}

\end{hinweis}

\subsection{Übersetzungsdateien}\label{uxfcbersetzungsdateien}

Übersetzungsdateien sind Mediendateien im Ordner \texttt{locale} des
jeweiligen Mods.

Luanti unterstützt derzeit zwei Formate: \emph{gettext} und ein eigenes,
älteres Format. Seit Luanti 5.10 gibt es gettext, ein verbreitetes und
bewährtes Format. Wenn möglich, solltest du gettext verwenden. Einziger
Nachteil: Luanti-Clients vor Version 5.10 können damit noch nicht
umgehen.

Bei gettext erwartet Luanti Übersetzungsdateien im Format
\texttt{[textdomain].\allowbreak{}[sprache].\allowbreak{}po} im Ordner
\texttt{locale}. Zum Erstellen und Pflegen kannst du externe Werkzeuge
verwenden. Die Luanti-Dokumentation enthält einen Artikel über
Weblate\footnote{\url{https://docs.luanti.org/for-creators/translation/weblate/}}
(englisch).

Lies außerdem die offizielle Luanti-Dokumentation zu
Übersetzungen\footnote{\url{https://api.luanti.org/translations/}}
(englisch).

Um Strings automatisch aus deinem Quellcode zu ziehen und
\texttt{.\allowbreak{}pot}-Vorlagen zu erzeugen, kannst du
\texttt{gettext} installieren und diesen Befehl anpassen:

\begin{Shaded}
\begin{Highlighting}[]
\NormalTok{xgettext {-}L lua {-}kS {-}kPS:1,2 \textbackslash{}}
\NormalTok{  {-}kcore.translate:1c,2 {-}kcore.translate\_n:1c,2,3 \textbackslash{}}
\NormalTok{  {-}d your\_textdomain {-}o locale/your\_textdomain.pot *.lua}
\end{Highlighting}
\end{Shaded}

Für jeden verwendeten Funktionsnamen kannst du ein weiteres \texttt{-k}
angeben. Die Angaben nach dem Funktionsnamen betreffen Pluralformen;
Einzelheiten stehen im Handbuch von gettext.

\subsubsection{Altes Format}\label{altes-format}

Übersetzungsdateien im alten Format heißen
\texttt{[textdomain].\allowbreak{}[sprache].\allowbreak{}tr}.

Sie beginnen mit einem Kommentar, der die Textdomain angibt; die
weiteren Zeilen bilden den englischen Ausgangstext auf die Übersetzung
ab.

Zum Beispiel \texttt{mymod.\allowbreak{}de.\allowbreak{}tr}:

\begin{Shaded}
\begin{Highlighting}[]
\NormalTok{\# textdomain: mymod}
\NormalTok{Hello everyone!=Hallo zusammen!}
\NormalTok{I like grapefruit=Ich mag Grapefruit}
\end{Highlighting}
\end{Shaded}

Übersetzungsdateien solltest du aus dem Quellcode deines Mods oder
Spiels erzeugen, etwa mit update\_translations\footnote{\url{https://github.com/minetest-tools/update_translations}}.
Das Werkzeug sucht in deinem Lua-Code nach \texttt{S(} und erstellt
automatisch eine Vorlage, die Übersetzer in ihre Sprache übertragen
können. Es aktualisiert die Übersetzungsdateien auch, wenn sich dein
Quellcode ändert.

Schreib in \texttt{S()} immer wörtlichen Text (in
\texttt{\textquotedbl{}}) statt einer Variablen; so finden Werkzeuge die
zu übersetzenden Texte.

\section{Platzhalter}\label{platzhalter}

Oft müssen variable Informationen in einen übersetzbaren Text. Dabei
darf der Text nicht einfach verkettet werden, denn sonst können
Übersetzer die Reihenfolge der Variablen im Satz nicht ändern. Verwende
stattdessen das Platzhaltersystem der Übersetzung:

\begin{Shaded}
\begin{Highlighting}[]
\VariableTok{core}\OperatorTok{.}\NormalTok{register\_on\_joinplayer}\OperatorTok{(}\KeywordTok{function}\OperatorTok{(}\VariableTok{player}\OperatorTok{)}
    \VariableTok{core}\OperatorTok{.}\NormalTok{chat\_send\_all}\OperatorTok{(}
\NormalTok{        S}\OperatorTok{(}
            \StringTok{"Everyone, say hi to @1!"}\OperatorTok{,}
            \VariableTok{player}\OperatorTok{:}\NormalTok{get\_player\_name}\OperatorTok{()}
        \OperatorTok{)}
    \OperatorTok{)}
\KeywordTok{end}\OperatorTok{)}
\end{Highlighting}
\end{Shaded}

Ein wörtliches \texttt{@} schreibst du in einer Übersetzung als
\texttt{@@}.

Innerhalb eines Satzes solltest du nicht verketten, mehrere Sätze
verkettest du aber am besten. Das hilft Übersetzern, weil die Strings
kürzer bleiben.

\begin{Shaded}
\begin{Highlighting}[]
\NormalTok{S}\OperatorTok{(}\StringTok{"Hello @1!"}\OperatorTok{,} \VariableTok{player\_name}\OperatorTok{)}
    \OperatorTok{..} \StringTok{" "}
    \OperatorTok{..}\NormalTok{ S}\OperatorTok{(}\StringTok{"You have @1 new messages."}\OperatorTok{,} \OperatorTok{\#}\VariableTok{msgs}\OperatorTok{),}
\end{Highlighting}
\end{Shaded}

\section{Gute Praxis und verbreitete
Irrtümer}\label{gute-praxis-und-verbreitete-irrtuxfcmer}

\begin{itemize}
\tightlist
\item
  Verkette keinen Text, sondern verwende Platzhalter. So haben
  Übersetzer volle Kontrolle über die Reihenfolge.
\item
  Variablen verändern oft den umgebenden Text, etwa durch Geschlecht und
  Pluralbildung. Das ist schwer zu handhaben und wird deshalb oft
  übergangen oder mit geschlechtsneutralen Formulierungen umgangen.
\item
  Übersetzungen können viel länger oder kürzer sein als der englische
  Text. Lass genug Platz.
\item
  Andere Sprachen schreiben Zahlen womöglich anders, etwa mit Komma als
  Dezimaltrennzeichen: \texttt{1.\allowbreak{}000,23},
  \texttt{1'000'000,32}.
\item
  Geh nicht davon aus, dass andere Sprachen Groß- und Kleinschreibung
  genauso verwenden wie Englisch.
\end{itemize}

\section{Übersetzung auf dem
Server}\label{uxfcbersetzung-auf-dem-server}

Manchmal brauchst du die Übersetzung eines Textes auf dem Server, etwa
um Text zu sortieren oder zu durchsuchen. Die Sprache eines Spielers
erhältst du mit
\texttt{get\_\allowbreak{}player\_\allowbreak{}information}, und mit
\texttt{get\_\allowbreak{}translated\_\allowbreak{}string} übersetzt du
markierten Text.

\begin{Shaded}
\begin{Highlighting}[]
\KeywordTok{local} \VariableTok{list} \OperatorTok{=} \OperatorTok{\{}
\NormalTok{    S}\OperatorTok{(}\StringTok{"Hello world!"}\OperatorTok{),}
\NormalTok{    S}\OperatorTok{(}\StringTok{"Potato"}\OperatorTok{),}
\OperatorTok{\}}

\VariableTok{core}\OperatorTok{.}\NormalTok{register\_chatcommand}\OperatorTok{(}\StringTok{"find"}\OperatorTok{,} \OperatorTok{\{}
    \VariableTok{func} \OperatorTok{=} \KeywordTok{function}\OperatorTok{(}\VariableTok{name}\OperatorTok{,} \VariableTok{param}\OperatorTok{)}
        \KeywordTok{local} \VariableTok{info} \OperatorTok{=} \VariableTok{core}\OperatorTok{.}\NormalTok{get\_player\_information}\OperatorTok{(}\VariableTok{name}\OperatorTok{)}
        \KeywordTok{local} \VariableTok{language} \OperatorTok{=} \VariableTok{info} \KeywordTok{and} \VariableTok{info}\OperatorTok{.}\VariableTok{language} \KeywordTok{or} \StringTok{"en"}

        \ControlFlowTok{for} \ConstantTok{\_}\OperatorTok{,} \VariableTok{line} \KeywordTok{in} \FunctionTok{ipairs}\OperatorTok{(}\VariableTok{list}\OperatorTok{)} \ControlFlowTok{do}
            \KeywordTok{local} \VariableTok{trans} \OperatorTok{=}
                \VariableTok{core}\OperatorTok{.}\NormalTok{get\_translated\_string}\OperatorTok{(}\VariableTok{language}\OperatorTok{,} \VariableTok{line}\OperatorTok{)}
            \ControlFlowTok{if} \VariableTok{trans}\OperatorTok{:}\FunctionTok{find}\OperatorTok{(}\VariableTok{param}\OperatorTok{,} \DecValTok{1}\OperatorTok{,} \KeywordTok{true}\OperatorTok{)} \ControlFlowTok{then}
                \ControlFlowTok{return} \KeywordTok{true}\OperatorTok{,} \VariableTok{line}
            \ControlFlowTok{end}
        \ControlFlowTok{end}
        \ControlFlowTok{return} \KeywordTok{false}\OperatorTok{,}\NormalTok{ S}\OperatorTok{(}\StringTok{"Nothing found."}\OperatorTok{)}
    \KeywordTok{end}\OperatorTok{,}
\OperatorTok{\})}
\end{Highlighting}
\end{Shaded}

\section{Fazit}\label{fazit}

Mit der Übersetzungs-API machst du Mods und Spiele mehr Menschen
zugänglich. Sie richtig einzusetzen, erfordert aber etwas Sorgfalt.
